% Generated draft: requires mathematical and source review.
\documentclass[11pt]{article}
\usepackage[margin=25mm]{geometry}
\usepackage{fontspec}
\setmainfont{Latin Modern Roman}
\usepackage{amsmath,amssymb,mathtools,unicode-math}
\setmathfont{Latin Modern Math}
\usepackage{enumitem,needspace,xurl,hyperref}
\setlength{\parindent}{0pt}
\setlength{\parskip}{5pt}
\setlength{\emergencystretch}{3em}
\newcommand{\N}{\mathbb{N}}
\newcommand{\Z}{\mathbb{Z}}
\newcommand{\Q}{\mathbb{Q}}
\newcommand{\R}{\mathbb{R}}
\newcommand{\C}{\mathbb{C}}
\newcommand{\F}{\mathbb{F}}
\newcommand{\A}{\mathbb{A}}
\newcommand{\PP}{\mathbb{P}}
\newcommand{\E}{\mathbb{E}}
\newcommand{\eps}{\varepsilon}
\newcommand{\dd}{\,\mathrm{d}}
\begin{document}
\section*{UnsolvedMath ID: 30005070}
\textbf{Source problem code:} catalogue entry
\subsection{Definitions and mathematical statement}
Let \(Q_T=(0,T)\times\R^2\), where \(T>0\) denotes the time horizon, and write \(\|\cdot\|_{L^p(Q_T)}\) for the Lebesgue norm in time and space. The supplied record denotes the forcing term by \(f\), the solution of a perturbed Kolmogorov equation by \(w\), and the second derivative in the \(x\)-direction by \(D_x^2w\).

The requested task is to give a purely analytic proof of the regularity estimate
\[
\|D_x^2w\|_{L^p(Q_T)}\leq M_p\|f\|_{L^p(Q_T)},\qquad M_p>0,
\]
for the perturbed Kolmogorov equation and the admissible data, solutions, exponents, and time horizons of the cited source. The reconstructed statement additionally requests, ideally, the constant \(M_p=C_p\), described there as sharp.

\textbf{Unresolved hypotheses and scope.} The supplied excerpt refers to equations (1), (2), and (3), but supplies none of those equations. In particular, it does not specify the perturbed operator, the admissible perturbations, initial conditions, the solution class, the range of \(p\), or the permitted dependence of \(M_p\) on \(T\) and the perturbation. The meaning and normalization of \(C_p\) are also absent. Thus the display records the requested estimate, but does not constitute a fully quantified theorem for a specified equation. No particular Kolmogorov operator or missing hypothesis is imposed here.

The original excerpt describes the stability result for equation (3) as having two parts, but supplies only part (i), concerning existence of \(M_p\). It also mentions related Schauder estimates, estimates of other \(L^p\) types, and anisotropic Schauder estimates for equation (1), without stating them. Its broader request for an analytic proof of ``our regularity results'' may therefore exceed the displayed estimate. The dimension-independent constants explicitly mentioned in the excerpt concern earlier heat-equation estimates; the supplied text does not specify a corresponding family of perturbed Kolmogorov equations in varying dimension.
\subsection{Short English statement}
Find a purely analytic proof of the supplied second-derivative \(L^p\) estimate for a perturbed Kolmogorov equation, ideally retaining the source's comparison constant \(C_p\). The equation, essential hypotheses, and definition of \(C_p\) are missing, so the intended theorem requires source clarification. The record reports a probabilistic proof and labels the analytic-proof request open; these are supplied assessments.
\subsection{Sources}
\begin{itemize}
\item Local UnsolvedMath registry, supplied record with permanent ID \(30005070\), source problem code \texttt{catalogue entry}, in the Oberwolfach Reports Open Problems collection. No individual proposer is supplied; the proposed year is 2022.
\item Supplied source citation: \emph{MFO-RIMS Tandem Workshop: Nonlocality in Analysis, Probability and Statistics (2023)}, Oberwolfach Reports, pp.~709--715, DOI \url{https://doi.org/10.4171/owr/2022/15}. The supplied literature-source label is \emph{Oberwolfach Report 15/2022}. The discrepancy between the citation's 2023 and the report's 2022 designation remains unresolved.
\item Supplied literature source: Marino--Menozzi--Priola, \emph{Poisson process and sharp constants}, \url{https://doi.org/10.1016/j.spa.2021.06.009}. The excerpt attributes the probabilistic stability argument to its reference [12] and the introduction of the Poisson-process approach to [9]; their full bibliographic entries are not supplied.
\item Literature assessment dated \texttt{2026-08-22}: the supplied curation labels the problem open, reports that a probabilistic representation proves the stated estimates, and says that a purely analytic proof remained unavailable and no later primary proof was found. This is supplied dated curation, not independently verified literature assessment.
\item The registry labels the reconstructed statement \texttt{corrected\_verified} and reports that it was checked against the official Oberwolfach report. These are supplied verification claims; the references have not been independently consulted for this draft.
\end{itemize}
\end{document}
